\chapter{Dos objetivos e responsabilidades}

\section{Dos objetivos}

\begin{enumerate}

\item
    Estabelecer diretrizes que permitam a todas as pessoas usuárias dos
    sistemas e ferramentas de tecnologia da informação e comunicação (TIC)
    do CRP--06 seguir padrões de comportamento relacionados à segurança da
    informação adequados às necessidades institucionais e à proteção legal
    das/os envolvidas/os;
\item
    normatizar e orientar todas as pessoas usuárias para a proteção dos
    sistemas, documentos, operações e comunicação institucional no que
    tange aos processos de trabalho e uso das redes, principalmente com a
    prevenção de responsabilidade legal;
\item
    definir normas e procedimentos específicos de segurança da informação,
    bem como a implementação de controles e processos para seu
    atendimento.
\item
    ofertar orientações sobre a utilização dos recursos computacionais, de
    telecomunicação e de infraestrutura de TI;
\item
    garantir a integridade da informação para que seja mantida em seu
    estado original, na guarda ou transmissão, contra alterações
    indevidas, intencionais ou acidentais;
\item
    garantir a confidencialidade e que o acesso à informação seja obtido
    somente por pessoas autorizadas e exclusivamente em função do
    interesse público;
\item
    garantir a disponibilidade para que as/os usuárias/os autorizadas/os
    obtenham acesso à informação e aos ativos correspondentes sempre que
    necessário;
\item
    assegurar a aplicabilidade das referências e diretrizes da governança
    da informação preconizadas pelo Tribunal de Contas da União (TCU);
\item
    organizar fluxos institucionais para tramitação de processos,
    informações, solicitações e respostas;
\item
    garantir a lisura, transparência e responsabilidade para com os
    processos institucionais.
\end{enumerate}

\section{Das responsabilidades}

Os recursos institucionais disponíveis na autarquia são de uso exclusivo
da e para as atividades e ações institucionais.

\begin{enumerate}

\item
    O CRP-06 não se responsabiliza pelas perdas causadas pelo uso
    indevido, negligente ou imprudente dos recursos e/ou serviços
    concedidos a suas/seus usuárias/os;
\end{enumerate}

\begin{enumerate}
\item
    o CRP-06, por meio da Unidade de Tecnologia da Informação e
    Comunicação, poderá registrar todo o uso dos sistemas e serviços,
    visando garantir a disponibilidade e a segurança das informações
    utilizadas, inclusive revisando e atualizando estas normas sempre que
    algum fato e/ou evento relevante acontecer;
\item
    todos os contratos firmados com as/os usuárias/os do CRP-06 deverão
    ser previamente formalizados via termo de acordo do uso e
    confidencialidade (TAUC), condição imprescindível para que possa ser
    concedido o acesso aos ativos de informação disponibilizados pela
    instituição;
\item
    qualquer incidente que afete a segurança da informação deverá ser
    comunicado inicialmente à coordenação de TIC e ela, após análise
    encaminhará à Gerência Administrativa e de Tecnologia da Informação
    (GATI) para que as medidas corretivas possam ser tomadas. Se
    necessário, a Unidade de Gestão de Pessoas e\emph{/}ou a diretoria
    poderão ser envolvidas para apoiar a decisão;
\item
    a fim de garantir maior segurança no uso dos sistemas da autarquia,
    faz-se necessária a segregação dos ambientes de desenvolvimento (uso
    da TI), testes (uso da TI e colaboradores-chave para testes),
    homologação (uso da TI) e produção (uso institucional);
\item
    o CRP-06 reserva-se o direito de analisar dados e evidências para
    obtenção de provas a serem utilizadas nos processos investigatórios,
    bem como adotar as medidas legais cabíveis;
\item
    é vedado o uso de todo e qualquer meio não institucional de contas e
    dispositivos pessoais como \emph{e-mails} que não sejam
    institucionais, WhatsApp, Skype, redes sociais e outros fora do
    domínio da rede formal;
\item
    é importante destacar que nenhuma pessoa deve ser lesada no uso de
    seus recursos pessoais, pois a autarquia oferece ferramentas
    necessárias para a comunicação institucional;
\item
    qualquer outro instrumento institucional (\emph{e-mail}, \emph{chat},
    aplicativo de mensagem instantânea, entre outros) não possui a
    prerrogativa deliberativa nem substitui, sob qualquer hipótese,
    aqueles já previstos, como ofício, memorando, despacho, portaria,
    instrução normativa ou resolução;
\item
    nas situações de não cumprimento, parcial ou integral, por parte da/o
    usuária/o, dos requisitos previstos nesta Política e das orientações
    sobre a segurança da informação, bem como das regras internas da
    autarquia, a/o usuária/o será submetida/o às medidas administrativas e
    legais cabíveis;
\item
    o uso pessoal dos recursos institucionais em hipótese alguma possui
    caráter deliberativo ou legal e é expressamente vedado e, caso ocorra,
    resultará na aplicação de medidas previstas na legislação pertinente
    para apuração dos fatos e responsabilização das/os envolvidas/os;
\item
    será de inteira responsabilidade de cada usuária/o o prejuízo ou dano
    que vier a sofrer ou causar ao CRP-06 e/ou a terceiros, em decorrência
    do não cumprimento às diretrizes e normas aqui referidas.
\end{enumerate}

\subsection{Das responsabilidades
específicas}

\subsubsection{Das/os conselheiras/os, gestoras/es, funcionárias/os,
colaboradoras/es, terceirizadas/os, estagiárias/os, prestadoras/es de
serviços e demais
convidadas/os}

\begin{enumerate}

\item
    Zelar pelo cumprimento desta Política, de forma a garantir a segurança
    da informação e assegurar o uso adequado dos meios, realizando
    \emph{backups} periódicos dentro da infraestrutura disponível, não
    compartilhando senhas e não disseminando notícias ou informações
    inconsistentes, sigilosas ou institucionais fora dos instrumentos da
    autarquia;
\end{enumerate}

\begin{enumerate}

\setcounter{enumi}{21}
\item
    exigir das/os usuárias/os sob sua responsabilidade o devido
    cumprimento desta Política;
\item
    exigir e condicionar o uso à assinatura do termo de uso e
    confidencialidade e o seguimento das normas estabelecidas, bem como ao
    compromisso de manter sigilo e confidencialidade sobre todos os ativos
    de informações do CRP-06, elucidando dúvidas e encaminhando para
    suporte as dificuldades técnicas, quando necessário.
\end{enumerate}

\subsubsection{Da Unidade Administrativa de Tecnologia da
Informação}

\begin{enumerate}

\item
    Informar às/aos gestores os serviços específicos que serão prestados e
    os procedimentos de resposta aos incidentes, atualizando, assim, o
    manual de serviços sempre que necessário;
\end{enumerate}

\begin{enumerate}

%\setcounter{enumi}{23}
\item
    configurar os equipamentos, ferramentas e sistemas concedidos às/os
    usuárias/os com todos os controles necessários para cumprir e
    assegurar os requisitos de segurança e uso estabelecidos por esta
    Política;
\item
    permitir, quando necessária, a execução de atividades operacionais sob
    sua responsabilidade, como, por exemplo, manutenção de computadores,
    realização de cópias de segurança, criação de contas para acesso
    institucional ou testes no ambiente;
\item
    seguir as normas e/ou orientações do Sistema Conselhos de Psicologia e
    da legislação vigente, no que diz respeito à segurança da informação
    para todos os Sistemas e aplicações com acesso público;
\item
    planejar, implantar, fornecer e monitorar a capacidade de armazenagem,
    processamento e transmissão necessária para garantir a segurança
    requerida pelas áreas de atuação, fundamentando tecnicamente seus
    procedimentos e encaminhamentos;
\item
    atribuir cada conta ou dispositivo de acesso (computadores, celulares,
    \emph{tablets} e outros), dispositivos dos sistemas, bases de dados ou
    qualquer outro ativo de informação a uma/um responsável identificável
    como pessoa física, de acordo com a Lei Geral de Proteção de Dados
    (LGPD);
\item
    proteger continuamente todos os ativos de informação do CRP-06 contra
    código malicioso, garantindo que os novos ativos só entrem para o
    ambiente de produção após estarem livres de código malicioso e/ou
    indesejado e após auditados por esta unidade administrativa;
\item
    propor a definição de regras formais para instalação de
    \emph{software} e \emph{hardware} em ambiente de produção
    institucional, implementando as diretrizes do Código de ética da
    Sociedade Brasileira de Computação em consonância com o Código de
    Ética Profissional da/o Psicóloga/o;
\item
    realizar manutenções corretiva, evolutiva e preventiva mensalmente, a
    fim de revisar tecnicamente os dispositivos tecnológicos de trabalho e
    mitigar possíveis riscos para assegurar o bom funcionamento dos
    recursos, a partir de plano de trabalho e cronograma devidamente
    aprovados pelo plenário;
\item
    auxiliar na instalação e configuração de assinaturas digitais e
    certificados digitais, por meio de ferramentas específicas;
\item
    promover os meios e estabelecer as diretrizes de orientação para as/os
    usuárias/os quanto à política de \emph{backup} de dados, e auditá-los
    quando necessário;
\item
    garantir, de maneira rápida e eficaz, após solicitação formal, o
    bloqueio de acesso de usuária/o por motivo de desligamento, incidente,
    investigação ou outra situação que exija medida restritiva para fins
    de resguardo dos ativos do CRP-06;
\item
    dar suporte operacional atitudinal na implementação, manutenção e
    promoção da cultura e desenvolvimento das ações e processos para uso
    de todos os sistemas, redes e ativos da autarquia, com base nos
    princípios éticos da área, de modo coerente com esta política.
\end{enumerate}

\paragraph{Da responsabilidade de garantia da segurança de
informação e
comunicação}

\begin{enumerate}

\item
    Publicar e promover as edições necessárias desta Política de Segurança
    em Tecnologia da Informação e Comunicação, com base em princípios
    técnicos;
\item
    promover a conscientização das/os usuárias/os em relação à relevância
    da segurança da informação, mediante campanhas e informativos, gestão
    de termos de cessão de uso e responsabilização;
\item
    analisar criticamente incidentes recorrentes e outros problemas em
    conjunto com usuária/o, fornecedores e, no caso dos Sistemas Federais,
    junto ao Comitê de TI do Conselho Federal de Psicologia;
\item
    analisar técnica e criticamente incidentes recorrentes e outros
    problemas de fornecedoras/es prestadoras/es de serviços em conjunto
    com as gerências e coordenações de cada unidade, elaborar parecer
    técnico fundamentado para dar subsídio à gerência de competência e dar
    ciência ao plenário, efetuando a gestão de contratos de uso e
    prestação de serviços;
\item
    exercer a continua práxis de alinhamento com as diretrizes
    institucionais da autarquia e normas regulamentadoras de segurança da
    informação e comunicação a partir de leis e normativas dos órgãos de
    controle.
\end{enumerate}